\section{Inleiding}

Het eerste deel van de cursus gaat over de macroscopische beschrijving van de thermische fysica. Dit wil zeggen op grote schaal. Later vanaf hoofdstuk 10 zullen we ook de microscopische beschrijvingen behandelen.
\subsection{Microscopisch en Macroscopisch}

\begin{itemize}
    \item Macroscopisch: 
    
    Behoort tot de klassieke thermodynamica. Hierin worden grootheden zoals druk, volume, temperatuur en energie gebruikt. Deze grootheden zijn gemiddeldes van microscopische grootheden.
    \item Microscopisch:
    
    Behoort tot de statistische mechanica. Hierin worden grootheden zoals de positie en snelheid van individuele deeltjes gebruikt, en ook de individuele eigen energien. Deze grootheden worden vaak achteerhaald met behulp van partitiefuncties en waarschijnlijkheidsverdelingen. 
\end{itemize}

\subsection{Toestanden}

Wanneer een coordinaat evenredig is met de massa, dan is ze \textbf{extensief}, de andere coordinaten zijn \textbf{intensief}. De toestand van een systeem wordt beschreven worden aan de hand van deze coördinaten.

Voorbeeld: Stel we hebben twee bekers water, waarvan de ene een volume van 1 liter heeft en de andere een volume van 2 liter op kamertemperatuur. In dit geval zijn de coördinaten van het systeem de druk, het volume en de temperatuur. 
Het volume is een extensieve grootheid, want wanneer we de 2 bekers samenvoegen, is het nieuwe volume gewoonweg de som van de volumes van de twee bekers. De druk en temperatuur zijn daarentegen intensieve grootheden, want de druk en temperatuur blijven op dat moment hetzelfde als die van de individuele bekers.

Het deel van de ruimte waarin we geïnteresseerd zijn, noemen we het \textbf{systeem}. De rest van de ruimte noemen we de \textbf{omgeving}. Het systeem en de omgeving vormen samen het \textbf{universum}.

\begin{definition}
    \textbf{Geïsoleerd systeem}: Een systeem dat geen energie of stof kan uitwisselen met de omgeving. (geen enkele interactie zoals thermische, mechanische of chemische met de omgeving mogelijk)
\end{definition}

\begin{definition}
    \textbf{Gesloten systeem}: Een systeem dat energie kan uitwisselen met de omgeving, maar geen stof kan uitwisselen.
\end{definition}

\begin{definition}
    \textbf{Diathermische wand}: Een wand die warmte kan doorlaten, maar geen materie. (laat alleen thermiche interacties toe)
\end{definition}

\begin{definition}
    \textbf{Adiabatische wand}: Een wand die geen warmte kan doorlaten (geen thermische interacties), maar de wand kan wel verplaatst worden (laat alleen mechanische interacties toe), de wand kan ook materie doorlaten, maar is niet per se nodig.
\end{definition}

Verschillende soorten interacties:
\begin{itemize}
    \item Mechanische interacties (vaste wanden \(\implies\) geen mechanische enrgie-uitwisseling, beweegbare wanden \(\implies\) wel mechanische energie-uitwisseling)
    \item Thermische interacties (adiabatische wanden \(\implies\) geen thermische interacties, diathermische wanden \(\implies\) wel thermische interacties)
    \item Chemische interacties (wanden die materie kunnen doorlaten \(\implies\) chemische interacties, wanden die geen materie kunnen doorlaten \(\implies\) geen chemische interacties)
\end{itemize}

\begin{definition}
    \textbf{Thermodynamisch evenwicht}: Systeem  als teglijkertijd mechanisch, thermisch en chemisch evenwicht is.
\end{definition}

\begin{definition}
  \textbf{Thermodynamische processen}: Veranderingen in de toestand van een systeem. Er zijn verschillende soorten processen, zoals isotherm (temperatuur blijft constant), isobaar (druk blijft constant), isochorisch (volume blijft constant) en adiabatisch (geen warmte-uitwisseling).  
\end{definition}

Als processen langzaam genoeg verlopen, kunnen we aannemen dat het systeem zich in een evenwichtstoestand bevindt tijdens het hele proces. We noemen dit een \textbf{quasi-statisch proces}.

De tijd om op evenwicht te komen, noemen we de \textbf{relaxatietijd} \(t_R\). (Dit komt niet meer aan bod verder dan hier)


